\chapter{Appendices}
\label{appendix}

\section{Per-Epoch Training Loss}
\label{sec:loss-full}

Table~\ref{tab:loss} in Chapter~\ref{ch3} reports the training loss at five checkpoints.
Table~\ref{tab:loss-full} gives the complete record for the six runs under seed 0, so that the shape
of the curves discussed in Section~\ref{sec:convergence} can be inspected directly. The other two
seeds follow the same trajectory and are not reproduced; their final values appear in
Table~\ref{tab:fid} and the underlying record is available in the accompanying repository.

The epoch at which one network overtakes another is not itself stable across seeds, and
Table~\ref{tab:crossover} records where it falls. The proposed network is overtaken by its
shortcut-free variant somewhere between the third and the sixth epoch on both datasets, which is
consistent with the early advantage being real but short-lived. The reference network is more
erratic: on CIFAR-10 it takes the lead as early as the second epoch under one seed and as late as
the twelfth under another, although it finishes ahead in every case. These figures are a further
reminder that quantities read from a single run should not be quoted to a precision the experiment
does not support.

\begin{table}[htbp]
    \centering
    \caption{The epoch at which one network first attains a lower training loss than another.}
    \label{tab:crossover}
    \small
    \begin{tabular}{llccc}
        \toprule
        Comparison & Dataset & Seed 0 & Seed 1 & Seed 2 \\
        \midrule
        plain overtakes ResUNet & CIFAR-10 & 4 & 4 & 3 \\
        plain overtakes ResUNet & CelebA   & 5 & 3 & 6 \\
        \midrule
        DDPM UNet leads both    & CIFAR-10 & 11 & 2 & 12 \\
        DDPM UNet leads both    & CelebA   & 2 & 1 & 2 \\
        \bottomrule
    \end{tabular}
\end{table}

\begin{table}[htbp]
    \centering
    \caption{Epoch-averaged noise-prediction loss for the six runs under seed 0. Lower is better.}
    \label{tab:loss-full}
    \footnotesize
    \begin{tabular}{ccccccc}
        \toprule
        & \multicolumn{3}{c}{CIFAR-10} & \multicolumn{3}{c}{CelebA} \\
        \cmidrule(lr){2-4} \cmidrule(lr){5-7}
        Epoch & ResUNet & plain & DDPM UNet & ResUNet & plain & DDPM UNet \\
        \midrule
        1 & 0.13775 & 0.14633 & 0.18508 & 0.14205 & 0.16413 & 0.15588 \\
        2 & 0.08237 & 0.08327 & 0.08805 & 0.06335 & 0.06774 & 0.06071 \\
        3 & 0.07665 & 0.07676 & 0.07861 & 0.05415 & 0.05513 & 0.04980 \\
        4 & 0.07322 & 0.07265 & 0.07399 & 0.04957 & 0.04963 & 0.04503 \\
        5 & 0.07111 & 0.07033 & 0.07135 & 0.04735 & 0.04705 & 0.04316 \\
        6 & 0.07017 & 0.06936 & 0.07014 & 0.04535 & 0.04496 & 0.04136 \\
        7 & 0.06790 & 0.06719 & 0.06799 & 0.04430 & 0.04379 & 0.04057 \\
        8 & 0.06705 & 0.06637 & 0.06672 & 0.04283 & 0.04234 & 0.03927 \\
        9 & 0.06655 & 0.06582 & 0.06591 & 0.04207 & 0.04159 & 0.03863 \\
        10 & 0.06619 & 0.06555 & 0.06593 & 0.04110 & 0.04055 & 0.03779 \\
        11 & 0.06570 & 0.06511 & 0.06491 & 0.04094 & 0.04044 & 0.03761 \\
        12 & 0.06517 & 0.06455 & 0.06448 & 0.03977 & 0.03924 & 0.03665 \\
        13 & 0.06458 & 0.06409 & 0.06386 & 0.03994 & 0.03948 & 0.03699 \\
        14 & 0.06424 & 0.06365 & 0.06360 & 0.03984 & 0.03935 & 0.03677 \\
        15 & 0.06424 & 0.06373 & 0.06322 & 0.03915 & 0.03866 & 0.03639 \\
        16 & 0.06352 & 0.06308 & 0.06281 & 0.03938 & 0.03891 & 0.03647 \\
        17 & 0.06315 & 0.06275 & 0.06206 & 0.03904 & 0.03855 & 0.03610 \\
        18 & 0.06305 & 0.06264 & 0.06202 & 0.03822 & 0.03775 & 0.03549 \\
        19 & 0.06278 & 0.06236 & 0.06196 & 0.03821 & 0.03773 & 0.03546 \\
        20 & 0.06249 & 0.06212 & 0.06165 & 0.03844 & 0.03799 & 0.03573 \\
        21 & 0.06250 & 0.06213 & 0.06158 & 0.03818 & 0.03772 & 0.03561 \\
        22 & 0.06265 & 0.06235 & 0.06160 & 0.03771 & 0.03727 & 0.03505 \\
        23 & 0.06170 & 0.06142 & 0.06062 & 0.03749 & 0.03712 & 0.03502 \\
        24 & 0.06204 & 0.06165 & 0.06080 & 0.03741 & 0.03700 & 0.03489 \\
        25 & 0.06218 & 0.06187 & 0.06091 & 0.03775 & 0.03726 & 0.03519 \\
        26 & 0.06132 & 0.06108 & 0.06040 & 0.03740 & 0.03704 & 0.03495 \\
        27 & 0.06156 & 0.06128 & 0.06075 & 0.03700 & 0.03659 & 0.03452 \\
        28 & 0.06190 & 0.06161 & 0.06071 & 0.03668 & 0.03626 & 0.03422 \\
        29 & 0.06165 & 0.06140 & 0.06059 & 0.03741 & 0.03700 & 0.03496 \\
        30 & 0.06108 & 0.06084 & 0.05967 & 0.03716 & 0.03679 & 0.03476 \\
        \bottomrule
    \end{tabular}
\end{table}

\section{Network Specifications}
\label{sec:specs}

The comparison in this work rests on the claim that the proposed and plain networks are identical in
size and that the reference network is comparable. Tables~\ref{tab:spec-resunet} and
\ref{tab:spec-ddpm} give the breakdown from which that claim can be checked.

The proposed and shortcut-free networks share every table entry, because an identity shortcut
introduces no parameters. Only the presence of the addition $x + h$ at the end of each block
differs, which is why the two rows of Table~\ref{tab:params} are equal to the digit.

\begin{table}[htbp]
    \centering
    \caption{Residual UNet, by stage. The figures apply unchanged to the shortcut-free variant.
    Each decoder concatenates the encoder features of the corresponding stage, so its block operates
    on the sum of the upsampled and skipped channel counts.}
    \label{tab:spec-resunet}
    \small
    \begin{tabular}{lccr}
        \toprule
        Stage & Channels & Resolution & Parameters \\
        \midrule
        \multicolumn{4}{l}{\emph{CIFAR-10, two resolution stages}} \\
        Timestep embedding MLP & --- & --- & 33\,024 \\
        Encoder 1 & $3 \rightarrow 64$ & $32 \times 32$ & 84\,160 \\
        Encoder 2 & $64 \rightarrow 128$ & $16 \times 16$ & 386\,048 \\
        Bridge & $128$ & $8 \times 8$ & 312\,192 \\
        Decoder 2 & $\rightarrow 256$ & $16 \times 16$ & 1\,279\,872 \\
        Decoder 1 & $\rightarrow 128$ & $32 \times 32$ & 377\,792 \\
        Final $1 \times 1$ convolution & $\rightarrow 3$ & $32 \times 32$ & 387 \\
        \emph{Total} & & & \emph{2\,473\,475} \\
        \midrule
        \multicolumn{4}{l}{\emph{CelebA, three resolution stages}} \\
        Timestep embedding MLP & --- & --- & 33\,024 \\
        Encoder 1 & $3 \rightarrow 32$ & $64 \times 64$ & 23\,648 \\
        Encoder 2 & $32 \rightarrow 64$ & $32 \times 32$ & 100\,864 \\
        Encoder 3 & $64 \rightarrow 128$ & $16 \times 16$ & 386\,048 \\
        Bridge & $128$ & $8 \times 8$ & 312\,192 \\
        Decoder 3 & $\rightarrow 256$ & $16 \times 16$ & 1\,279\,872 \\
        Decoder 2 & $\rightarrow 128$ & $32 \times 32$ & 377\,792 \\
        Decoder 1 & $\rightarrow 64$ & $64 \times 64$ & 98\,784 \\
        Final $1 \times 1$ convolution & $\rightarrow 3$ & $64 \times 64$ & 195 \\
        \emph{Total} & & & \emph{2\,612\,419} \\
        \bottomrule
    \end{tabular}
\end{table}

Two features of Table~\ref{tab:spec-resunet} are worth noting. The first is that the coarsest
decoder dominates the budget, holding $52\%$ of the parameters on CIFAR-10 and $49\%$ on CelebA,
because concatenation doubles the channel count before the block that follows it. The second is that
adding a third resolution stage for CelebA increases the total by only $5.6\%$: the channel widths
were narrowed from $(64, 128)$ to $(32, 64, 128)$ so that the deeper network would remain comparable
in size to the shallower one, which is what permits the two datasets to be discussed on the same
footing.

\begin{table}[htbp]
    \centering
    \caption{Reference DDPM UNet, by functional group. The same configuration is used for both
    datasets, since none of its parameter counts depends on the input resolution.}
    \label{tab:spec-ddpm}
    \small
    \begin{tabular}{lr}
        \toprule
        Group & Parameters \\
        \midrule
        Timestep embedding MLP & 131\,712 \\
        Input convolution & 896 \\
        Downward residual blocks (8) & 523\,328 \\
        Downward attention & 33\,536 \\
        Downsampling convolutions (3) & 83\,104 \\
        Middle block (residual, attention, residual) & 181\,504 \\
        Upward residual blocks (12) & 1\,240\,640 \\
        Upward attention & 50\,304 \\
        Upsampling convolutions (3) & 110\,784 \\
        Output normalisation and convolution & 931 \\
        \midrule
        \emph{Total} & \emph{2\,356\,739} \\
        \bottomrule
    \end{tabular}
\end{table}

The reference network reaches a similar total by a different distribution. Its upward path carries
$53\%$ of the parameters, because a skip connection is taken from every residual block rather than
once per resolution stage, so each of the twelve upward blocks receives a concatenated input. The
attention layers account for only $3.6\%$ of the weights while contributing a large share of the
runtime, which is why this network costs roughly $2.3$ times as much per epoch despite being the
smallest of the three.

\section{Memorisation Check: Remaining Configurations}
\label{sec:memo-full}

Figure~\ref{fig:memorisation} in Chapter~\ref{ch3} shows the closest generated and training pairs
for the proposed network on CelebA. The remaining five configurations are given here, and support
the same conclusion: in no case is a generated sample a reproduction of a training image.

\begin{figure}[htbp]
    \centering
    \includegraphics[width=\linewidth]{memo_cifar10_resunet.png}\\[2mm]
    \includegraphics[width=\linewidth]{memo_cifar10_plain.png}\\[2mm]
    \includegraphics[width=\linewidth]{memo_cifar10_ddpm.png}
    \caption{CIFAR-10: the eight samples closest to the training set (top row of each pair) with
    their nearest training image (bottom row), for the proposed, shortcut-free and reference
    networks in turn.}
    \label{fig:memo-cifar}
\end{figure}

\begin{figure}[htbp]
    \centering
    \includegraphics[width=\linewidth]{memo_celeba_plain.png}\\[3mm]
    \includegraphics[width=\linewidth]{memo_celeba_ddpm.png}
    \caption{CelebA: the shortcut-free and reference networks. No pair depicts the same person.}
    \label{fig:memo-celeba}
\end{figure}

\section{Computing Environment}
\label{sec:environment}

All eighteen runs were carried out on the machine described in Table~\ref{tab:environment}. Six
configurations trained for thirty epochs under each of three seeds amount to approximately five
hours of training, and the evaluation adds a further two and a half, dominated by drawing the two
thousand samples per model required for the Fr\'echet Inception Distance. The complete experiment
therefore occupies about seven and a half hours on a single consumer graphics card, which is what
makes the repetition recommended in the Conclusions practical rather than aspirational.

\begin{table}[htbp]
    \centering
    \caption{Hardware and software used for every run reported in this work.}
    \label{tab:environment}
    \small
    \begin{tabular}{ll}
        \toprule
        Component & Version \\
        \midrule
        GPU & NVIDIA GeForce RTX 4070, 12\,GB \\
        Driver & 591.86 \\
        Operating system & Windows 11 \\
        Python & 3.12.10 \\
        PyTorch & 2.14.0+cu130 \\
        torchvision & 0.29.0+cu130 \\
        CUDA runtime & 13.0 \\
        cuDNN & 9.24.0 \\
        torchmetrics & 1.9.0 \\
        torch-fidelity & 0.4.0 \\
        scikit-image & 0.26.0 \\
        NumPy & 2.5.2 \\
        \bottomrule
    \end{tabular}
\end{table}

Training used bfloat16 mixed precision with the channels-last memory format, and TF32 matrix
multiplication with the cuDNN autotuner enabled. These settings affect only the arithmetic used to
evaluate a fixed model, not the model itself, and were applied identically to all six runs.

The implementation is a single Jupyter notebook containing the network definitions, the diffusion
process, the training loop and every evaluation reported here. It includes a self-check routine that
verifies the noise schedule is monotone and converges to the standard normal distribution, that the
closed-form forward process is consistent, that the network output depends on the timestep, that the
sampling path executes, and that the parameter budgets of the three architectures satisfy the
constraints stated in Section~\ref{sec:simulation}. Datasets are decoded once into unsigned
eight-bit tensors and cached, so that only the first execution pays the preprocessing cost, and
training is skipped for any architecture whose checkpoint already records the full number of epochs.
